\textbf{A continuación se presenta la descomposición del crecimiento de cada una de las trece ramas comparables de actividad económica CIIU.} En particular, se presenta el PIB real de la actividad, el número de ocupados, el PIB por trabajador, las horas semanales por trabajador y el PIB por hora trabajada para los años 2010 y 2025. La lectura conjunta de estas variables permite distinguir si los aumentos en la producción responden principalmente a variaciones en el empleo, a cambios en las horas trabajadas o un aumento de la productividad por hora.

%Las trece ramas comparables parten de las agrupaciones reportadas por el DANE en cuentas nacionales. La diferencia frente a las doce agrupaciones usuales es que este informe separa la agrupación R+S+T en dos ramas: arte, entretenimiento, recreación y otros servicios, por un lado, y hogares como empleadores, por el otro. Esta separación es posible porque tanto las cuentas nacionales como la GEIH permiten construir numeradores y denominadores laborales comparables para ambas ramas. Como las series se expresan en pesos constantes, pueden presentarse pequeñas diferencias de aditividad entre niveles de agregación.

\subsection{Agropecuario}

\textbf{La actividad agropecuaria comprende cultivos, apoyo agropecuario y actividades mixtas; café; ganadería; silvicultura; y pesca y acuicultura.} En 2025 la actividad agropecuaria representó 6,3\% del PIB, con 63,9 billones de pesos de 2015, y concentró 13,6\% de los ocupados, con 3,1 millones de personas. El PIB por trabajador fue 20,4 millones de pesos de 2015 por ocupado, por debajo del agregado (44,5 millones). Las horas semanales por trabajador fueron 42,0, muy cerca del agregado (43,7). El PIB por hora trabajada fue 9,4 mil pesos de 2015 por hora, por debajo del agregado (19,6 mil).

\begin{table}[H]
\centering
\caption{Agropecuario: PIB, ocupados, horas y productividad laboral, 2010--2025}
\label{tab:sector_a_productividad}
\small
\begin{tabular}{lrrr}
\toprule
Indicador & 2010 & 2025 & Crec. anual \\
\midrule
PIB real  (Billones de pesos de 2015) & 39,9 & 63,9 & 3,18\% \\
Ocupados (Millones) & 3,0 & 3,1 & 0,38\% \\
PIB por trabajador (Millones de pesos de 2015) & 13,5 & 20,4 & 2,79\% \\
Horas semanales por trabajador & 45,0 & 42,0 & -0,46\% \\
PIB por hora trabajada (Miles de pesos de 2015) & 5,8 & 9,4 & 3,27\% \\
\bottomrule
\end{tabular}
\end{table}

\textbf{El crecimiento de la productividad del sector agropecuario fue superior al crecimiento de la productividad agregada.} Entre 2010 y 2025 el PIB de la actividad agropecuaria registró un crecimiento anualizado del 3,18\%, muy cerca del agregado (3,16\%). Los ocupados en la actividad registraron un crecimiento anual del 0,38\%, muy por debajo del agregado (2,09\%). El PIB por trabajador registró un crecimiento anual del 2,79\%, muy por encima del crecimiento del PIB por trabajador de toda la economía (1,05\%). Las horas semanales por trabajador registraron una caída anual del 0,46\%, muy cerca de la caída agregada (-0,45\%). El PIB por hora trabajada registró un crecimiento anual del 3,27\%, muy por encima del crecimiento del PIB por hora trabajada de toda la economía (1,51\%).

\begin{table}[H]
\centering
\caption{Productividad laboral de las actividades agropecuarias, 2010--2025}
\label{tab:sector_a_zoom61}
\small
\begin{tabular}{p{0.40\textwidth}rrrr}
\toprule
Actividad económica & PIB/trab. 2025 & Crec. & PIB/hora 2025 & Crec. \\
\midrule
Silvicultura & 70,2 & 5,15\% & 32,9 & 6,27\% \\
Pesca y acuicultura & 24,6 & 4,79\% & 10,9 & 5,75\% \\
Ganadería & 27,1 & 4,45\% & 12,0 & 5,10\% \\
Cultivos, apoyo agropecuario y mixtos & 20,5 & 2,24\% & 9,6 & 2,61\% \\
Café & 9,3 & 1,75\% & 4,1 & 1,98\% \\
\bottomrule
\end{tabular}
\caption*{\footnotesize Nota: PIB por trabajador en millones de pesos constantes de 2015; PIB por hora en miles de pesos constantes de 2015. %Cuando la GEIH no separa ocupados y horas al mismo nivel de las cuentas nacionales, se agrupan las subactividades del DANE hasta el nivel laboral comparable. 
Fuente: cálculos propios con DANE y GEIH.}
\end{table}

\textbf{El balance de la apertura del sector agropecuario es favorable: todas las subactividades superan el crecimiento agregado tanto en PIB por trabajador como en PIB por hora trabajada.} En esta apertura, el mayor crecimiento del PIB por trabajador se observa en Silvicultura (5,15\%), muy por encima del crecimiento agregado del PIB por trabajador (1,05\%); el menor se registra en Café (1,75\%), por encima del crecimiento agregado del PIB por trabajador (1,05\%). Por hora trabajada, el mejor desempeño corresponde a Silvicultura (6,27\%), muy por encima del crecimiento agregado del PIB por hora trabajada (1,51\%); el más rezagado corresponde a Café (1,98\%), por encima del crecimiento agregado del PIB por hora trabajada (1,51\%).

\subsection{Minas}

\textbf{La actividad de minas y canteras comprende petróleo, gas y apoyo conexo; otras minas y canteras; carbón; y minerales metalíferos.} En 2025 la actividad de minas y canteras representó 3,4\% del PIB, con 34,7 billones de pesos de 2015, y concentró 1,3\% de los ocupados, con 0,3 millones de personas. El PIB por trabajador fue 115,8 millones de pesos de 2015 por ocupado, por encima del agregado (44,5 millones). Las horas semanales por trabajador fueron 46,0, por encima del agregado (43,7). El PIB por hora trabajada fue 48,4 mil pesos de 2015 por hora, por encima del agregado (19,6 mil).

\begin{table}[H]
\centering
\caption{Minas: PIB, ocupados, horas y productividad laboral, 2010--2025}
\label{tab:sector_b_productividad}
\small
\begin{tabular}{lrrr}
\toprule
Indicador & 2010 & 2025 & Crec. anual \\
\midrule
PIB real  (Billones de pesos de 2015) & 38,4 & 34,7 & -0,67\% \\
Ocupados (Millones) & 0,2 & 0,3 & 3,04\% \\
PIB por trabajador (Millones de pesos de 2015) & 200,8 & 115,8 & -3,60\% \\
Horas semanales por trabajador & 49,7 & 46,0 & -0,52\% \\
PIB por hora trabajada (Miles de pesos de 2015) & 77,6 & 48,4 & -3,09\% \\
\bottomrule
\end{tabular}
\end{table}

\textbf{La productividad de minas y canteras cayó entre 2010 y 2025.} Entre 2010 y 2025 el PIB de la actividad de minas y canteras registró una caída anualizada del 0,67\%, muy por debajo del agregado (3,16\%). Los ocupados en la actividad registraron un crecimiento anual del 3,04\%, por encima del agregado (2,09\%). El PIB por trabajador registró una caída anual del 3,60\%, muy por debajo del crecimiento del PIB por trabajador de toda la economía (1,05\%). Las horas semanales por trabajador registraron una caída anual del 0,52\%, muy cerca de la caída agregada (-0,45\%). El PIB por hora trabajada registró una caída anual del 3,09\%, muy por debajo del crecimiento del PIB por hora trabajada de toda la economía (1,51\%).

\begin{table}[H]
\centering
\caption{Productividad laboral de las actividades mineras, 2010--2025}
\label{tab:sector_b_zoom61}
\small
\begin{tabular}{p{0.40\textwidth}rrrr}
\toprule
Actividad económica & PIB/trab. 2025 & Crec. & PIB/hora 2025 & Crec. \\
\midrule
Petróleo, gas y apoyo conexo & 576,1 & 0,07\% & 231,6 & 0,59\% \\
Otras minas y canteras & 45,6 & -0,89\% & 19,7 & -0,45\% \\
Carbón & 103,5 & -5,88\% & 41,8 & -5,33\% \\
Minerales metalíferos & 17,6 & -6,44\% & 7,5 & -5,98\% \\
\bottomrule
\end{tabular}
\caption*{\footnotesize Nota: PIB por trabajador en millones de pesos constantes de 2015; PIB por hora en miles de pesos constantes de 2015. %Cuando la GEIH no separa ocupados y horas al mismo nivel de las cuentas nacionales, se agrupan las subactividades del DANE hasta el nivel laboral comparable.
Fuente: cálculos propios con DANE y GEIH.}
\end{table}

\textbf{El balance de la apertura de la actividad de minas y canteras es negativo: ninguna subactividad supera el crecimiento agregado y tres de las cuatro subactividades caen tanto en PIB por trabajador como en PIB por hora trabajada.} En esta apertura, el mayor crecimiento del PIB por trabajador se observa en Petróleo, gas y apoyo conexo (0,07\%), por debajo del crecimiento agregado del PIB por trabajador (1,05\%); el menor se registra en Minerales metalíferos (-6,44\%), muy por debajo del crecimiento agregado del PIB por trabajador (1,05\%). Por hora trabajada, el mejor desempeño corresponde a Petróleo, gas y apoyo conexo (0,59\%), por debajo del crecimiento agregado del PIB por hora trabajada (1,51\%); el más rezagado corresponde a Minerales metalíferos (-5,98\%), muy por debajo del crecimiento agregado del PIB por hora trabajada (1,51\%).

\subsection{Manufactura}

\textbf{La actividad manufacturera agrupa alimentos, bebidas y tabaco; textiles, confecciones y cuero; madera, papel e impresión; refinación, químicos y minerales; metalurgia, maquinaria y equipo; y muebles y otras manufactureras.} En 2025 la actividad manufacturera representó 11,1\% del PIB, con 113,1 billones de pesos de 2015, y concentró 10,9\% de los ocupados, con 2,5 millones de personas. El PIB por trabajador fue 45,0 millones de pesos de 2015 por ocupado, muy cerca del agregado (44,5 millones). Las horas semanales por trabajador fueron 44,3, muy cerca del agregado (43,7). El PIB por hora trabajada fue 19,5 mil pesos de 2015 por hora, muy cerca del agregado (19,6 mil).

\begin{table}[H]
\centering
\caption{Manufactura: PIB, ocupados, horas y productividad laboral, 2010--2025}
\label{tab:sector_c_productividad}
\small
\begin{tabular}{lrrr}
\toprule
Indicador & 2010 & 2025 & Crec. anual \\
\midrule
PIB real  (Billones de pesos de 2015) & 88,0 & 113,1 & 1,69\% \\
Ocupados (Millones) & 2,2 & 2,5 & 0,91\% \\
PIB por trabajador (Millones de pesos de 2015) & 40,1 & 45,0 & 0,77\% \\
Horas semanales por trabajador & 46,7 & 44,3 & -0,35\% \\
PIB por hora trabajada (Miles de pesos de 2015) & 16,5 & 19,5 & 1,12\% \\
\bottomrule
\end{tabular}
\end{table}

\textbf{El crecimiento de la productividad de la manufactura fue inferior al crecimiento de la productividad agregada.} Entre 2010 y 2025 el PIB de la actividad manufacturera registró un crecimiento anualizado del 1,69\%, muy por debajo del agregado (3,16\%). Los ocupados en la actividad registraron un crecimiento anual del 0,91\%, muy por debajo del agregado (2,09\%). El PIB por trabajador registró un crecimiento anual del 0,77\%, por debajo del crecimiento del PIB por trabajador de toda la economía (1,05\%). Las horas semanales por trabajador registraron una caída anual del 0,35\%, muy cerca de la caída agregada (-0,45\%). El PIB por hora trabajada registró un crecimiento anual del 1,12\%, por debajo del crecimiento del PIB por hora trabajada de toda la economía (1,51\%).

\begin{table}[H]
\centering
\caption{Productividad laboral de las actividades manufactureras, 2010--2025}
\label{tab:sector_c_zoom25}
\small
\begin{tabular}{p{0.40\textwidth}rrrr}
\toprule
Actividad económica & PIB/trab. 2025 & Crec. & PIB/hora 2025 & Crec. \\
\midrule
Metalurgia, maquinaria y equipo & 48,8 & 1,69\% & 20,5 & 2,16\% \\
Alimentos, bebidas y tabaco & 50,8 & 1,60\% & 21,8 & 1,70\% \\
Madera, papel e impresión & 46,9 & 1,46\% & 21,8 & 2,21\% \\
Textiles, confecciones y cuero & 17,5 & 0,74\% & 7,8 & 1,19\% \\
Refinación, químicos y minerales & 106,6 & -0,75\% & 44,8 & -0,30\% \\
Muebles y otras manufactureras & 17,4 & -1,38\% & 7,6 & -1,13\% \\
\bottomrule
\end{tabular}
\caption*{\footnotesize Nota: PIB por trabajador en millones de pesos constantes de 2015; PIB por hora en miles de pesos constantes de 2015. Fuente: cálculos propios con DANE y GEIH.}
\end{table}

\textbf{El balance de la apertura de la manufactura es heterogéneo: tres de las seis subactividades superan el crecimiento agregado del PIB por trabajador y tres de las seis subactividades superan el crecimiento agregado del PIB por hora trabajada, mientras otras quedan rezagadas.} En esta apertura, el mayor crecimiento del PIB por trabajador se observa en Metalurgia, maquinaria y equipo (1,69\%), por encima del crecimiento agregado del PIB por trabajador (1,05\%); el menor se registra en Muebles y otras manufactureras (-1,38\%), muy por debajo del crecimiento agregado del PIB por trabajador (1,05\%). Por hora trabajada, el mejor desempeño corresponde a Madera, papel e impresión (2,21\%), por encima del crecimiento agregado del PIB por hora trabajada (1,51\%); el más rezagado corresponde a Muebles y otras manufactureras (-1,13\%), muy por debajo del crecimiento agregado del PIB por hora trabajada (1,51\%).

\subsection{Servicios públicos}

\textbf{La actividad de servicios públicos combina electricidad, gas y vapor con agua, saneamiento y desechos.} En 2025 la actividad de servicios públicos representó 3,0\% del PIB, con 30,3 billones de pesos de 2015, y concentró 1,3\% de los ocupados, con 0,3 millones de personas. El PIB por trabajador fue 99,5 millones de pesos de 2015 por ocupado, por encima del agregado (44,5 millones). Las horas semanales por trabajador fueron 44,0, muy cerca del agregado (43,7). El PIB por hora trabajada fue 43,5 mil pesos de 2015 por hora, por encima del agregado (19,6 mil).

\begin{table}[H]
\centering
\caption{Servicios públicos: PIB, ocupados, horas y productividad laboral, 2010--2025}
\label{tab:sector_d_e_productividad}
\small
\begin{tabular}{lrrr}
\toprule
Indicador & 2010 & 2025 & Crec. anual \\
\midrule
PIB real  (Billones de pesos de 2015) & 21,9 & 30,3 & 2,18\% \\
Ocupados (Millones) & 0,1 & 0,3 & 6,16\% \\
PIB por trabajador (Millones de pesos de 2015) & 176,5 & 99,5 & -3,75\% \\
Horas semanales por trabajador & 48,5 & 44,0 & -0,64\% \\
PIB por hora trabajada (Miles de pesos de 2015) & 70,0 & 43,5 & -3,13\% \\
\bottomrule
\end{tabular}
\end{table}

\textbf{La productividad de los servicios públicos cayó entre 2010 y 2025.} Entre 2010 y 2025 el PIB de la actividad de servicios públicos registró un crecimiento anualizado del 2,18\%, por debajo del agregado (3,16\%). Los ocupados en la actividad registraron un crecimiento anual del 6,16\%, muy por encima del agregado (2,09\%). El PIB por trabajador registró una caída anual del 3,75\%, muy por debajo del crecimiento del PIB por trabajador de toda la economía (1,05\%). Las horas semanales por trabajador registraron una caída anual del 0,64\%, muy cerca de la caída agregada (-0,45\%). El PIB por hora trabajada registró una caída anual del 3,13\%, muy por debajo del crecimiento del PIB por hora trabajada de toda la economía (1,51\%).

\begin{table}[H]
\centering
\caption{Productividad laboral de las actividades de servicios públicos, 2010--2025}
\label{tab:sector_d_e_zoom25}
\small
\begin{tabular}{p{0.40\textwidth}rrrr}
\toprule
Actividad económica & PIB/trab. 2025 & Crec. & PIB/hora 2025 & Crec. \\
\midrule
Electricidad, gas y vapor & 262,2 & -0,13\% & 109,6 & 0,28\% \\
Agua, saneamiento y desechos & 37,9 & -6,41\% & 16,9 & -5,74\% \\
\bottomrule
\end{tabular}
\caption*{\footnotesize Nota: PIB por trabajador en millones de pesos constantes de 2015; PIB por hora en miles de pesos constantes de 2015. Fuente: cálculos propios con DANE y GEIH.}
\end{table}

\textbf{El balance de la apertura de los servicios públicos es negativo: ninguna subactividad supera el crecimiento agregado, todas registran caídas en PIB por trabajador y una de las dos subactividades también cae por hora trabajada.} En esta apertura, el mayor crecimiento del PIB por trabajador se observa en Electricidad, gas y vapor (-0,13\%), muy por debajo del crecimiento agregado del PIB por trabajador (1,05\%); el menor se registra en Agua, saneamiento y desechos (-6,41\%), muy por debajo del crecimiento agregado del PIB por trabajador (1,05\%). Por hora trabajada, el mejor desempeño corresponde a Electricidad, gas y vapor (0,28\%), muy por debajo del crecimiento agregado del PIB por hora trabajada (1,51\%); el más rezagado corresponde a Agua, saneamiento y desechos (-5,74\%), muy por debajo del crecimiento agregado del PIB por hora trabajada (1,51\%).

\subsection{Construcción}

\textbf{La actividad de construcción comprende edificaciones, obras civiles y actividades especializadas de construcción.} En 2025 la actividad de construcción representó 4,0\% del PIB, con 41,3 billones de pesos de 2015, y concentró 6,9\% de los ocupados, con 1,6 millones de personas. El PIB por trabajador fue 26,1 millones de pesos de 2015 por ocupado, por debajo del agregado (44,5 millones). Las horas semanales por trabajador fueron 45,7, muy cerca del agregado (43,7). El PIB por hora trabajada fue 11,0 mil pesos de 2015 por hora, por debajo del agregado (19,6 mil).

\begin{table}[H]
\centering
\caption{Construcción: PIB, ocupados, horas y productividad laboral, 2010--2025}
\label{tab:sector_f_productividad}
\small
\begin{tabular}{lrrr}
\toprule
Indicador & 2010 & 2025 & Crec. anual \\
\midrule
PIB real  (Billones de pesos de 2015) & 40,0 & 41,3 & 0,21\% \\
Ocupados (Millones) & 1,0 & 1,6 & 3,39\% \\
PIB por trabajador (Millones de pesos de 2015) & 41,8 & 26,1 & -3,08\% \\
Horas semanales por trabajador & 48,9 & 45,7 & -0,46\% \\
PIB por hora trabajada (Miles de pesos de 2015) & 16,4 & 11,0 & -2,63\% \\
\bottomrule
\end{tabular}
\end{table}

\textbf{La productividad de la construcción cayó entre 2010 y 2025.} Entre 2010 y 2025 el PIB de la actividad de construcción registró un crecimiento anualizado del 0,21\%, muy por debajo del agregado (3,16\%). Los ocupados en la actividad registraron un crecimiento anual del 3,39\%, muy por encima del agregado (2,09\%). El PIB por trabajador registró una caída anual del 3,08\%, muy por debajo del crecimiento del PIB por trabajador de toda la economía (1,05\%). Las horas semanales por trabajador registraron una caída anual del 0,46\%, muy cerca de la caída agregada (-0,45\%). El PIB por hora trabajada registró una caída anual del 2,63\%, muy por debajo del crecimiento del PIB por hora trabajada de toda la economía (1,51\%).

\begin{table}[H]
\centering
\caption{Productividad laboral de las actividades de construcción, 2010--2025}
\label{tab:sector_f_zoom25}
\small
\begin{tabular}{p{0.40\textwidth}rrrr}
\toprule
Actividad económica & PIB/trab. 2025 & Crec. & PIB/hora 2025 & Crec. \\
\midrule
Obras civiles & 56,2 & -2,58\% & 23,1 & -2,01\% \\
Actividades especializadas de construcción & 26,7 & -3,27\% & 11,7 & -2,90\% \\
Edificaciones & 19,9 & -3,52\% & 8,3 & -3,08\% \\
\bottomrule
\end{tabular}
\caption*{\footnotesize Nota: PIB por trabajador en millones de pesos constantes de 2015; PIB por hora en miles de pesos constantes de 2015. Fuente: cálculos propios con DANE y GEIH.}
\end{table}

\textbf{El balance de la apertura de la construcción es claramente negativo: todas las subactividades registran caídas de productividad y ninguna alcanza el crecimiento agregado.} En esta apertura, el mayor crecimiento del PIB por trabajador se observa en Obras civiles (-2,58\%), muy por debajo del crecimiento agregado del PIB por trabajador (1,05\%); el menor se registra en Edificaciones (-3,52\%), muy por debajo del crecimiento agregado del PIB por trabajador (1,05\%). Por hora trabajada, el mejor desempeño corresponde a Obras civiles (-2,01\%), muy por debajo del crecimiento agregado del PIB por hora trabajada (1,51\%); el más rezagado corresponde a Edificaciones (-3,08\%), muy por debajo del crecimiento agregado del PIB por hora trabajada (1,51\%).

\subsection{Comercio, transporte y alojamiento}

\textbf{La agrupación de comercio, transporte y alojamiento combina comercio y reparación; transporte y almacenamiento; y alojamiento y servicios de comida.} En 2025 la agrupación de comercio, transporte y alojamiento representó 17,5\% del PIB, con 179,1 billones de pesos de 2015, y concentró 32,4\% de los ocupados, con 7,4 millones de personas. El PIB por trabajador fue 24,1 millones de pesos de 2015 por ocupado, por debajo del agregado (44,5 millones). Las horas semanales por trabajador fueron 45,8, muy cerca del agregado (43,7). El PIB por hora trabajada fue 10,1 mil pesos de 2015 por hora, por debajo del agregado (19,6 mil).

\begin{table}[H]
\centering
\caption{Comercio, transporte y alojamiento: PIB, ocupados, horas y productividad laboral, 2010--2025}
\label{tab:sector_g_h_i_productividad}
\small
\begin{tabular}{lrrr}
\toprule
Indicador & 2010 & 2025 & Crec. anual \\
\midrule
PIB real  (Billones de pesos de 2015) & 107,6 & 179,1 & 3,45\% \\
Ocupados (Millones) & 5,2 & 7,4 & 2,46\% \\
PIB por trabajador (Millones de pesos de 2015) & 20,9 & 24,1 & 0,97\% \\
Horas semanales por trabajador & 50,6 & 45,8 & -0,66\% \\
PIB por hora trabajada (Miles de pesos de 2015) & 7,9 & 10,1 & 1,64\% \\
\bottomrule
\end{tabular}
\end{table}

\textbf{La productividad de comercio, transporte y alojamiento evolucionó de forma parecida a la productividad agregada.} Entre 2010 y 2025 el PIB de la agrupación de comercio, transporte y alojamiento registró un crecimiento anualizado del 3,45\%, por encima del agregado (3,16\%). Los ocupados en la actividad registraron un crecimiento anual del 2,46\%, por encima del agregado (2,09\%). El PIB por trabajador registró un crecimiento anual del 0,97\%, muy cerca del crecimiento del PIB por trabajador de toda la economía (1,05\%). Las horas semanales por trabajador registraron una caída anual del 0,66\%, una caída más pronunciada que la del agregado (-0,45\%). El PIB por hora trabajada registró un crecimiento anual del 1,64\%, muy cerca del crecimiento del PIB por hora trabajada de toda la economía (1,51\%).

\begin{table}[H]
\centering
\caption{Productividad laboral de las actividades de comercio, transporte y alojamiento, 2010--2025}
\label{tab:sector_g_h_i_zoom25}
\small
\begin{tabular}{p{0.40\textwidth}rrrr}
\toprule
Actividad económica & PIB/trab. 2025 & Crec. & PIB/hora 2025 & Crec. \\
\midrule
Comercio y reparación & 22,8 & 2,28\% & 9,5 & 2,76\% \\
Transporte y almacenamiento & 27,8 & -0,18\% & 10,9 & 0,83\% \\
Alojamiento y comida & 23,4 & -1,08\% & 10,7 & -0,42\% \\
\bottomrule
\end{tabular}
\caption*{\footnotesize Nota: PIB por trabajador en millones de pesos constantes de 2015; PIB por hora en miles de pesos constantes de 2015. Fuente: cálculos propios con DANE y GEIH.}
\end{table}

\textbf{El balance de la apertura de la agrupación de comercio, transporte y alojamiento es heterogéneo: solo una de las tres subactividades supera en crecimiento de la productividad al resto de la economía.} En esta apertura, el mayor crecimiento del PIB por trabajador se observa en Comercio y reparación (2,28\%), muy por encima del crecimiento agregado del PIB por trabajador (1,05\%); el menor se registra en Alojamiento y comida (-1,08\%), muy por debajo del crecimiento agregado del PIB por trabajador (1,05\%). Por hora trabajada, el mejor desempeño corresponde a Comercio y reparación (2,76\%), muy por encima del crecimiento agregado del PIB por hora trabajada (1,51\%); el más rezagado corresponde a Alojamiento y comida (-0,42\%), muy por debajo del crecimiento agregado del PIB por hora trabajada (1,51\%).

\subsection{Información y comunicaciones}

\textbf{La actividad de información y comunicaciones reúne edición de contenidos audiovisuales, radio, televisión, telecomunicaciones, desarrollo de software, procesamiento de datos, y otros servicios de información.} En 2025 la actividad de información y comunicaciones representó 3,1\% del PIB, con 31,3 billones de pesos de 2015, y concentró 1,7\% de los ocupados, con 0,4 millones de personas. El PIB por trabajador fue 78,6 millones de pesos de 2015 por ocupado, por encima del agregado (44,5 millones). Las horas semanales por trabajador fueron 44,4, muy cerca del agregado (43,7). El PIB por hora trabajada fue 34,0 mil pesos de 2015 por hora, por encima del agregado (19,6 mil).

\begin{table}[H]
\centering
\caption{Información y comunicaciones: PIB, ocupados, horas y productividad laboral, 2010--2025}
\label{tab:sector_j_productividad}
\small
\begin{tabular}{lrrr}
\toprule
Indicador & 2010 & 2025 & Crec. anual \\
\midrule
PIB real  (Billones de pesos de 2015) & 18,3 & 31,3 & 3,66\% \\
Ocupados (Millones) & 0,4 & 0,4 & -0,07\% \\
PIB por trabajador (Millones de pesos de 2015) & 45,4 & 78,6 & 3,73\% \\
Horas semanales por trabajador & 49,6 & 44,4 & -0,73\% \\
PIB por hora trabajada (Miles de pesos de 2015) & 17,6 & 34,0 & 4,49\% \\
\bottomrule
\end{tabular}
\end{table}

\textbf{El crecimiento de la productividad de información y comunicaciones fue superior al crecimiento de la productividad agregada.} Entre 2010 y 2025 el PIB de la actividad de información y comunicaciones registró un crecimiento anualizado del 3,66\%, por encima del agregado (3,16\%). Los ocupados en la actividad registraron una caída anual del 0,07\%, muy por debajo del agregado (2,09\%). El PIB por trabajador registró un crecimiento anual del 3,73\%, muy por encima del crecimiento del PIB por trabajador de toda la economía (1,05\%). Las horas semanales por trabajador registraron una caída anual del 0,73\%, una caída más pronunciada que la del agregado (-0,45\%). El PIB por hora trabajada registró un crecimiento anual del 4,49\%, muy por encima del crecimiento del PIB por hora trabajada de toda la economía (1,51\%).

\subsection{Financieras}

\textbf{La actividad financiera reúne intermediación financiera y banca, seguros, fondos de pensiones y servicios auxiliares como administración de mercados, corretaje y otras actividades de apoyo financiero.} En 2025 la actividad financiera representó 5,2\% del PIB, con 53,1 billones de pesos de 2015, y concentró 1,9\% de los ocupados, con 0,4 millones de personas. El PIB por trabajador fue 124,5 millones de pesos de 2015 por ocupado, por encima del agregado (44,5 millones). Las horas semanales por trabajador fueron 44,2, muy cerca del agregado (43,7). El PIB por hora trabajada fue 54,2 mil pesos de 2015 por hora, por encima del agregado (19,6 mil).

\begin{table}[H]
\centering
\caption{Financieras: PIB, ocupados, horas y productividad laboral, 2010--2025}
\label{tab:sector_k_productividad}
\small
\begin{tabular}{lrrr}
\toprule
Indicador & 2010 & 2025 & Crec. anual \\
\midrule
PIB real  (Billones de pesos de 2015) & 22,3 & 53,1 & 5,95\% \\
Ocupados (Millones) & 0,2 & 0,4 & 4,49\% \\
PIB por trabajador (Millones de pesos de 2015) & 101,0 & 124,5 & 1,40\% \\
Horas semanales por trabajador & 45,2 & 44,2 & -0,15\% \\
PIB por hora trabajada (Miles de pesos de 2015) & 43,0 & 54,2 & 1,56\% \\
\bottomrule
\end{tabular}
\end{table}

\textbf{La evolución de la productividad de las actividades financieras fue parecida a la del resto de la economía.} Entre 2010 y 2025 el PIB de la actividad financiera registró un crecimiento anualizado del 5,95\%, muy por encima del agregado (3,16\%). Los ocupados en la actividad registraron un crecimiento anual del 4,49\%, muy por encima del agregado (2,09\%). El PIB por trabajador registró un crecimiento anual del 1,40\%, por encima del crecimiento del PIB por trabajador de toda la economía (1,05\%). Las horas semanales por trabajador registraron una caída anual del 0,15\%, una caída menos pronunciada que la del agregado (-0,45\%). El PIB por hora trabajada registró un crecimiento anual del 1,56\%, muy cerca del crecimiento del PIB por hora trabajada de toda la economía (1,51\%).

\subsection{Inmobiliarias}

\textbf{La actividad inmobiliaria reúne alquiler, compra, venta, administración e intermediación de bienes inmuebles, tanto propios o arrendados como de terceros.} En 2025 la actividad inmobiliaria representó 8,8\% del PIB, con 90,1 billones de pesos de 2015, y concentró 1,4\% de los ocupados, con 0,3 millones de personas. El PIB por trabajador fue 284,6 millones de pesos de 2015 por ocupado, por encima del agregado (44,5 millones). Las horas semanales por trabajador fueron 45,8, muy cerca del agregado (43,7). El PIB por hora trabajada fue 119,5 mil pesos de 2015 por hora, por encima del agregado (19,6 mil).

\begin{table}[H]
\centering
\caption{Inmobiliarias: PIB, ocupados, horas y productividad laboral, 2010--2025}
\label{tab:sector_l_productividad}
\small
\begin{tabular}{lrrr}
\toprule
Indicador & 2010 & 2025 & Crec. anual \\
\midrule
PIB real  (Billones de pesos de 2015) & 59,9 & 90,1 & 2,75\% \\
Ocupados (Millones) & 0,2 & 0,3 & 2,64\% \\
PIB por trabajador (Millones de pesos de 2015) & 280,0 & 284,6 & 0,11\% \\
Horas semanales por trabajador & 51,1 & 45,8 & -0,73\% \\
PIB por hora trabajada (Miles de pesos de 2015) & 105,3 & 119,5 & 0,85\% \\
\bottomrule
\end{tabular}
\end{table}

\textbf{El crecimiento de la productividad de las actividades inmobiliarias fue inferior al crecimiento de la productividad agregada.} Entre 2010 y 2025 el PIB de la actividad inmobiliaria registró un crecimiento anualizado del 2,75\%, por debajo del agregado (3,16\%). Los ocupados en la actividad registraron un crecimiento anual del 2,64\%, por encima del agregado (2,09\%). El PIB por trabajador registró un crecimiento anual del 0,11\%, por debajo del crecimiento del PIB por trabajador de toda la economía (1,05\%). Las horas semanales por trabajador registraron una caída anual del 0,73\%, una caída más pronunciada que la del agregado (-0,45\%). El PIB por hora trabajada registró un crecimiento anual del 0,85\%, por debajo del crecimiento del PIB por hora trabajada de toda la economía (1,51\%).

\subsection{Profesionales y administrativas}

\textbf{La agrupación de actividades profesionales y administrativas combina actividades profesionales, científicas y técnicas con servicios administrativos y de apoyo.} En 2025 la agrupación de actividades profesionales y administrativas representó 6,9\% del PIB, con 70,1 billones de pesos de 2015, y concentró 7,2\% de los ocupados, con 1,7 millones de personas. El PIB por trabajador fue 42,4 millones de pesos de 2015 por ocupado, muy cerca del agregado (44,5 millones). Las horas semanales por trabajador fueron 39,6, por debajo del agregado (43,7). El PIB por hora trabajada fue 20,6 mil pesos de 2015 por hora, por encima del agregado (19,6 mil).

\begin{table}[H]
\centering
\caption{Profesionales y administrativas: PIB, ocupados, horas y productividad laboral, 2010--2025}
\label{tab:sector_m_n_productividad}
\small
\begin{tabular}{lrrr}
\toprule
Indicador & 2010 & 2025 & Crec. anual \\
\midrule
PIB real  (Billones de pesos de 2015) & 45,4 & 70,1 & 2,94\% \\
Ocupados (Millones) & 0,8 & 1,7 & 4,65\% \\
PIB por trabajador (Millones de pesos de 2015) & 54,3 & 42,4 & -1,63\% \\
Horas semanales por trabajador & 42,1 & 39,6 & -0,41\% \\
PIB por hora trabajada (Miles de pesos de 2015) & 24,8 & 20,6 & -1,22\% \\
\bottomrule
\end{tabular}
\end{table}

\textbf{La productividad de las actividades profesionales y administrativas cayó entre 2010 y 2025.} Entre 2010 y 2025 el PIB de la agrupación de actividades profesionales y administrativas registró un crecimiento anualizado del 2,94\%, por debajo del agregado (3,16\%). Los ocupados en la actividad registraron un crecimiento anual del 4,65\%, muy por encima del agregado (2,09\%). El PIB por trabajador registró una caída anual del 1,63\%, muy por debajo del crecimiento del PIB por trabajador de toda la economía (1,05\%). Las horas semanales por trabajador registraron una caída anual del 0,41\%, muy cerca de la caída agregada (-0,45\%). El PIB por hora trabajada registró una caída anual del 1,22\%, muy por debajo del crecimiento del PIB por hora trabajada de toda la economía (1,51\%).

\begin{table}[H]
\centering
\caption{Productividad laboral de las actividades profesionales y administrativas, 2010--2025}
\label{tab:sector_m_n_zoom61}
\small
\begin{tabular}{p{0.40\textwidth}rrrr}
\toprule
Actividad económica & PIB/trab. 2025 & Crec. & PIB/hora 2025 & Crec. \\
\midrule
Servicios administrativos y de apoyo & 36,0 & -1,32\% & 18,0 & -0,87\% \\
Profesionales, científicas y técnicas & 54,7 & -1,91\% & 25,2 & -1,60\% \\
\bottomrule
\end{tabular}
\caption*{\footnotesize Nota: PIB por trabajador en millones de pesos constantes de 2015; PIB por hora en miles de pesos constantes de 2015. %Cuando la GEIH no separa ocupados y horas al mismo nivel de las cuentas nacionales, se agrupan las subactividades del DANE hasta el nivel laboral comparable.
Fuente: cálculos propios con DANE y GEIH.}
\end{table}

\textbf{El balance de la apertura de las actividades profesionales y administrativas es claramente negativo: todas las subactividades registran caídas de productividad y ninguna alcanza el crecimiento agregado.} En esta apertura, el mayor crecimiento del PIB por trabajador se observa en Servicios administrativos y de apoyo (-1,32\%), muy por debajo del crecimiento agregado del PIB por trabajador (1,05\%); el menor se registra en Profesionales, científicas y técnicas (-1,91\%), muy por debajo del crecimiento agregado del PIB por trabajador (1,05\%). Por hora trabajada, el mejor desempeño corresponde a Servicios administrativos y de apoyo (-0,87\%), muy por debajo del crecimiento agregado del PIB por hora trabajada (1,51\%); el más rezagado corresponde a Profesionales, científicas y técnicas (-1,60\%), muy por debajo del crecimiento agregado del PIB por hora trabajada (1,51\%).

\subsection{Adm. pública, educación y salud}

\textbf{La agrupación de administración pública, educación y salud combina administración pública y defensa, educación, salud humana y servicios sociales.} En 2025 la agrupación de administración pública, educación y salud representó 16,2\% del PIB, con 165,3 billones de pesos de 2015, y concentró 12,3\% de los ocupados, con 2,8 millones de personas. El PIB por trabajador fue 58,4 millones de pesos de 2015 por ocupado, por encima del agregado (44,5 millones). Las horas semanales por trabajador fueron 43,0, muy cerca del agregado (43,7). El PIB por hora trabajada fue 26,1 mil pesos de 2015 por hora, por encima del agregado (19,6 mil).

\begin{table}[H]
\centering
\caption{Adm. pública, educación y salud: PIB, ocupados, horas y productividad laboral, 2010--2025}
\label{tab:sector_o_p_q_productividad}
\small
\begin{tabular}{lrrr}
\toprule
Indicador & 2010 & 2025 & Crec. anual \\
\midrule
PIB real  (Billones de pesos de 2015) & 85,4 & 165,3 & 4,51\% \\
Ocupados (Millones) & 1,9 & 2,8 & 2,54\% \\
PIB por trabajador (Millones de pesos de 2015) & 43,9 & 58,4 & 1,92\% \\
Horas semanales por trabajador & 43,3 & 43,0 & -0,03\% \\
PIB por hora trabajada (Miles de pesos de 2015) & 19,5 & 26,1 & 1,95\% \\
\bottomrule
\end{tabular}
\end{table}

\textbf{El crecimiento de la productividad de administración pública, educación y salud fue superior al crecimiento de la productividad agregada.} Entre 2010 y 2025 el PIB de la agrupación de administración pública, educación y salud registró un crecimiento anualizado del 4,51\%, muy por encima del agregado (3,16\%). Los ocupados en la actividad registraron un crecimiento anual del 2,54\%, por encima del agregado (2,09\%). El PIB por trabajador registró un crecimiento anual del 1,92\%, por encima del crecimiento del PIB por trabajador de toda la economía (1,05\%). Las horas semanales por trabajador registraron una caída anual del 0,03\%, una caída menos pronunciada que la del agregado (-0,45\%). El PIB por hora trabajada registró un crecimiento anual del 1,95\%, por encima del crecimiento del PIB por hora trabajada de toda la economía (1,51\%).

\begin{table}[H]
\centering
\caption{Productividad laboral de las actividades de administración pública, educación y salud, 2010--2025}
\label{tab:sector_o_p_q_zoom25}
\small
\begin{tabular}{p{0.40\textwidth}rrrr}
\toprule
Actividad económica & PIB/trab. 2025 & Crec. & PIB/hora 2025 & Crec. \\
\midrule
Salud y servicios sociales & 49,1 & 2,54\% & 21,3 & 2,47\% \\
Educación & 52,3 & 1,54\% & 25,8 & 1,26\% \\
Administración pública & 77,2 & 1,40\% & 32,2 & 2,06\% \\
\bottomrule
\end{tabular}
\caption*{\footnotesize Nota: PIB por trabajador en millones de pesos constantes de 2015; PIB por hora en miles de pesos constantes de 2015. Fuente: cálculos propios con DANE y GEIH.}
\end{table}

\textbf{El balance de la apertura de administración pública, educación y salud es favorable en PIB por trabajador, pero menos uniforme por hora trabajada: todas las subactividades superan el crecimiento agregado por trabajador y dos de las tres subactividades lo hacen por hora.} En esta apertura, el mayor crecimiento del PIB por trabajador se observa en Salud y servicios sociales (2,54\%), muy por encima del crecimiento agregado del PIB por trabajador (1,05\%); el menor se registra en Administración pública (1,40\%), por encima del crecimiento agregado del PIB por trabajador (1,05\%). Por hora trabajada, el mejor desempeño corresponde a Salud y servicios sociales (2,47\%), por encima del crecimiento agregado del PIB por hora trabajada (1,51\%); el más rezagado corresponde a Educación (1,26\%), por debajo del crecimiento agregado del PIB por hora trabajada (1,51\%).

\subsection{Arte, entretenimiento, recreación y otros servicios}

\textbf{La agrupación de arte, entretenimiento, recreación y otros servicios reúne actividades artísticas, de entretenimiento y recreación, juegos de azar, actividades deportivas, servicios personales, asociaciones y reparación de computadores y efectos personales.} En 2025 la agrupación de arte, entretenimiento, recreación y otros servicios representó 4,0\% del PIB, con 40,6 billones de pesos de 2015, y concentró 5,4\% de los ocupados, con 1,2 millones de personas. El PIB por trabajador fue 32,8 millones de pesos de 2015 por ocupado, por debajo del agregado (44,5 millones). Las horas semanales por trabajador fueron 39,6, por debajo del agregado (43,7). El PIB por hora trabajada fue 16,0 mil pesos de 2015 por hora, por debajo del agregado (19,6 mil).

\begin{table}[H]
\centering
\caption{Arte, entretenimiento, recreación y otros servicios: PIB, ocupados, horas y productividad laboral, 2010--2025}
\label{tab:sector_r_s_productividad}
\small
\begin{tabular}{lrrr}
\toprule
Indicador & 2010 & 2025 & Crec. anual \\
\midrule
PIB real  (Billones de pesos de 2015) & 11,0 & 40,6 & 9,09\% \\
Ocupados (Millones) & 1,0 & 1,2 & 1,65\% \\
PIB por trabajador (Millones de pesos de 2015) & 11,4 & 32,8 & 7,32\% \\
Horas semanales por trabajador & 39,3 & 39,6 & 0,04\% \\
PIB por hora trabajada (Miles de pesos de 2015) & 5,6 & 16,0 & 7,28\% \\
\bottomrule
\end{tabular}
\end{table}

\textbf{El crecimiento de la productividad de arte, entretenimiento, recreación y otros servicios fue superior al crecimiento de la productividad agregada.} Entre 2010 y 2025 el PIB de la agrupación de arte, entretenimiento, recreación y otros servicios registró un crecimiento anualizado del 9,09\%, muy por encima del agregado (3,16\%). Los ocupados en la actividad registraron un crecimiento anual del 1,65\%, por debajo del agregado (2,09\%). El PIB por trabajador registró un crecimiento anual del 7,32\%, muy por encima del crecimiento del PIB por trabajador de toda la economía (1,05\%). Las horas semanales por trabajador registraron un crecimiento anual del 0,04\%, por encima del agregado (-0,45\%). El PIB por hora trabajada registró un crecimiento anual del 7,28\%, muy por encima del crecimiento del PIB por hora trabajada de toda la economía (1,51\%).

\subsection{Hogares como empleadores}

\textbf{La actividad de hogares como empleadores reúne el trabajo doméstico remunerado en hogares y la producción no diferenciada de los hogares para uso propio.} En 2025 la actividad de hogares como empleadores representó 0,7\% del PIB, con 6,7 billones de pesos de 2015, y concentró 3,1\% de los ocupados, con 0,7 millones de personas. El PIB por trabajador fue 9,4 millones de pesos de 2015 por ocupado, por debajo del agregado (44,5 millones). Las horas semanales por trabajador fueron 40,5, por debajo del agregado (43,7). El PIB por hora trabajada fue 4,5 mil pesos de 2015 por hora, por debajo del agregado (19,6 mil).

\begin{table}[H]
\centering
\caption{Hogares como empleadores: PIB, ocupados, horas y productividad laboral, 2010--2025}
\label{tab:sector_t_productividad}
\small
\begin{tabular}{lrrr}
\toprule
Indicador & 2010 & 2025 & Crec. anual \\
\midrule
PIB real  (Billones de pesos de 2015) & 4,3 & 6,7 & 3,06\% \\
Ocupados (Millones) & 0,7 & 0,7 & 0,42\% \\
PIB por trabajador (Millones de pesos de 2015) & 6,4 & 9,4 & 2,63\% \\
Horas semanales por trabajador & 45,9 & 40,5 & -0,83\% \\
PIB por hora trabajada (Miles de pesos de 2015) & 2,7 & 4,5 & 3,49\% \\
\bottomrule
\end{tabular}
\end{table}

\textbf{El crecimiento de la productividad de hogares como empleadores fue superior al crecimiento de la productividad agregada.} Entre 2010 y 2025 el PIB de la actividad de hogares como empleadores registró un crecimiento anualizado del 3,06\%, muy cerca del agregado (3,16\%). Los ocupados en la actividad registraron un crecimiento anual del 0,42\%, muy por debajo del agregado (2,09\%). El PIB por trabajador registró un crecimiento anual del 2,63\%, muy por encima del crecimiento del PIB por trabajador de toda la economía (1,05\%). Las horas semanales por trabajador registraron una caída anual del 0,83\%, una caída más pronunciada que la del agregado (-0,45\%). El PIB por hora trabajada registró un crecimiento anual del 3,49\%, muy por encima del crecimiento del PIB por hora trabajada de toda la economía (1,51\%).

%{\footnotesize Nota general: para facilitar la lectura, el informe usa PIB por actividad económica; en sentido estricto, el numerador corresponde al valor agregado bruto de cada actividad reportado por el DANE. Valores en pesos constantes de 2015. Ocupados expandidos con el factor \texttt{fex}. Las horas semanales promedio se calculan como horas anuales totales divididas por ocupados y por 52 semanas. Se excluye 2020 por no contar con GEIH anual comparable en la base del proyecto. Fuente: cálculos propios con DANE y GEIH.}
